\documentclass[11pt]{article}

% parameters are fairly self-explanatory here, thankfully
\usepackage[width=7in,height=9.5in,top=0.75in,centering]{geometry}

% for math-ing
\usepackage{amsmath}
\usepackage{amssymb}
\usepackage{comment}

%%%%%% tables %%%%%%%%
\usepackage{tabularx}
\usepackage{array}
\renewcommand{\arraystretch}{1.5}
\setlength{\tabcolsep}{8pt}

%%%%%% plotting %%%%%%

\usepackage[dvipsnames]{xcolor}
\usepackage{tikz}
\usepackage{pgfplots}

\pgfplotsset{every axis/.append style={
                    axis x line=middle,    % put the x axis in the middle
                    axis y line=middle,    % put the y axis in the middle
                    axis line style={<->,color=black}, % arrows on the axis
                    xlabel={$x$},          % default put x on x-axis
                    ylabel={$y$},          % default put y on y-axis
            }}
            
 \pgfplotsset{compat=1.14} % sometimes, everything falls apart without this
 
 %\usepgfplotslibrary{fillbetween} %for shading between curves

%%%%%% commands %%%%%%

\def\todo{\textcolor{red}{TODO}}

% exam heading; course name is first argument, exam information is second argument
\newcommand{\Heading}[2]{
    \fbox{
        {\Large\textbf{#1}}
        }
    \hfill
    \fbox{
        \begin{minipage}{0.2\textwidth}
            \begin{center}
            \textbf{#2}
            \end{center}
        \end{minipage}
    }
    \par\vspace{0.1in}}

% for being lazy while beautifying integrals
\newcommand{\dint}{\displaystyle\int}

% for being lazy while beautifying limits
\newcommand{\dlim}{\displaystyle\lim}

%%%%%% stuff %%%%%%

\usepackage{fancyhdr}

%\setcounter{page}{0} % first page is page 0

\parindent0pt % no indent

\fboxsep=0.25cm % little space in fbox border

%%%%%% BEGIN! %%%%%%%

\begin{document}

\pagestyle{fancy}

% record goals here to create sense of transparency

\fancyhead[l]{%
  \ifcase\value{page}
  \or {\bf\color{ForestGreen} D3: Derivatives of Basic Functions}
  \or {\bf\color{ForestGreen} D2: Compute Derivatives}
  \or {\bf\color{ProcessBlue} F3: Summations}
  \or {\bf\color{Orange} P1: Discrete Random Variables}
  \or {\bf\color{Orange} P2: Continuous Random Variables}
  \or {\bf\color{Orange} P1: Discrete Random Variables}
  \or {\bf\color{Orange} P2: Continuous Random Variables}
  \or
  \else
  % Empty chead here!
  \fi
}

% \thispagestyle{empty} % don't display that this is page 0

\Heading{MATH 1210}{Final Review 2\\Spring 2026\\
Core: D3, D2\\
Non-core: F3, P1, P2\\
Section 930}

\vspace{0.1in}

\textbf{Name} \hrulefill


\begin{enumerate}
\item[\textbf{D3.}]
Write down the derivatives of the following functions.
No need to simplify.
\begin{center}
    \begin{tabular}{l c r l c r}
    %
        $\displaystyle\frac{d}{dx}\left[ -4x^{9} \right]$ & $=$ & \rule[-8pt]{1.5in}{0.4pt} &
        \vspace{2cm}
    %
        $\displaystyle\frac{d}{dx}\left[ \dfrac{2}{\sqrt{x^3}} \right]$ & $=$ & \rule[-8pt]{1.5in}{0.4pt}\\[0.3in]
        \vspace{2cm}
    %%%%%
        $\displaystyle\frac{d}{dx}\left[ 4^x \right]$ & $=$ & \rule[-8pt]{1.5in}{0.4pt} & 
    %%%%%
        $\displaystyle\frac{d}{dx}\left[ \frac{\pi^2}{6} \right]$ & $=$ & \rule[-8pt]{1.5in}{0.4pt}\\[0.3in]
        \vspace{2cm}
    %%%%%
        $\displaystyle\frac{d}{dx}\left[ x^2 + x+1 \right]$ & $=$ & \rule[-8pt]{1.5in}{0.4pt}&
    %
        $\displaystyle\frac{d}{dx}\left[ \ln(x) + x \right]$ & $=$ & \rule[-8pt]{1.5in}{0.4pt}\\[0.3in]
        \vspace{2cm}
    %%%%%
        $\displaystyle\frac{d}{dx}\left[ \log_2(x^5) \right]$ & $=$ & \rule[-8pt]{1.5in}{0.4pt} & 
    %
        $\displaystyle\frac{d}{dx}\left[ \frac{1}{x} + \frac{2}{x^2} \right]$ & $=$ & \rule[-8pt]{1.5in}{0.4pt}\\[0.3in]
        \vspace{2cm}

        $\displaystyle\frac{d}{dx}\left[ (3x^2-5x+1)^7 \right]$ & $=$ & \rule[-8pt]{1.5in}{0.4pt} & 
    %
        $\displaystyle\frac{d}{dx}\left[ \sqrt{2x^3+4x} \right]$ & $=$ & \rule[-8pt]{1.5in}{0.4pt}\\[0.3in]
    %%%%%
    \end{tabular}

\end{center}
\pagebreak
\item[\textbf{D3.}]Find the following derivatives:
\begin{itemize}
\item
$\dfrac{d}{dx} \left[ \left( \dfrac{x^3-6x^2}{2x+1} \right) \cdot \ln(3^{x+1}) \right]$
\vfill
\item 
$\dfrac{d}{dx} \left[ \dfrac{-7x^4+3x^3-2x}{x^2}\cdot \left(3^x-\dfrac{5}{x^4}\right) \right]$
\vfill
\end{itemize}
\newpage
\item[\textbf{F3.}]
\begin{enumerate}
    \item Find the value of the following expressions. \textbf{Show all work.} Simplify as much as possible so that your
    \textbf{final answer is a number}.

    \begin{enumerate}
        \item[(i)] $\displaystyle\sum_{i=4}^6\dfrac{2i}{i-3}$

        \vfill

        \item[(ii)] $\displaystyle\sum_{j=2}^54$ 

        \vfill
    \end{enumerate}

    \item Rewrite the following in summation notation. 

    \[e^2+3e^3+5e^4+7e^5\]

    \vfill
\end{enumerate}
\newpage
\item[\textbf{P1.}] Suppose you roll two 4-sided dice, and then compute the difference of the second roll from the first. Let \( X \) be the random variable taking this difference as its value.
    \begin{enumerate}

        \item Is \( X \) continuous or discrete? If it is continuous, what is its support? If it is discrete, list all its possible values.
            \vspace{3cm}

        \item What is the sample space for the experiment?
            \vspace{3cm}

        \item If it is discrete, find the PMF of \( X \).
        If it is continuous, find the PDF of $X$.
            \vspace{5cm}

        \item Find the expectation of \( X \).

    \end{enumerate}

    \pagebreak

\item[\textbf{P2.}] The expected value (or mean) of a continuous random variable $X$ is $\displaystyle \mu=E[X]=\int_1^5k \,dx$ where $k$ is an unknown constant that we will find later. 
    \begin{enumerate}

        \item What is the probability density function $p(x)$? Your answer will involve $k$.
            \vspace{3cm}

        \item Find the value of $k$ to ensure that $p(x)$ is a pdf. What is the support of \( X \)?
            \vspace{3cm}

        \item Find the cumulative distribution function (CDF) $C(x)$.
            \vspace{3cm}

        \item Use $C(x)$ to find $P(X \leq 4)$.
            \vspace{3cm}

        \item Find the expectation and variance (without \( k \)) of \( X \).

    \end{enumerate}
\newpage

%%%%%%%%%%%%%%%%%%%%%%%%%%%%%%%%%%%%%%%%%%%%%%%%%%%%%
%%%%%%%%%%% BEGIN PROBLEM P1 %%%%%%%%%%%%%%%%%%%%%%%%
%%%%%%%%%%%%%%%%%%%%%%%%%%%%%%%%%%%%%%%%%%%%%%%%%%%%%

\item[\textbf{P1.}]
Each lottery ticket has a 1\% chance of paying out \$100 and a 10\% chance of paying out \$10. Otherwise the payout is \$0.
Let $X$ be a discrete random variable representing the amount a lottery ticket pays out.

\begin{enumerate}
    \item What are the possible values of $X$?

    \vfill

    \item Write down the probability mass function $p(x)$ of $X$ using either a table or a piecewise function.

    \vfill
    \vfill

    \item Find the expectation of $X$.

    \vfill
    \vfill

    \item Find the variance of $X$.

    \vfill
    \vfill

\end{enumerate}

%%%%%%%%%%%%%%%%%%%%%%%%%%%%%%%%%%%%%%%%%%%%%%%%%%%%%
%%%%%%%%%%% END PROBLEM P1 %%%%%%%%%%%%%%%%%%%%%%%%%%
%%%%%%%%%%%%%%%%%%%%%%%%%%%%%%%%%%%%%%%%%%%%%%%%%%%%%

\newpage

%%%%%%%%%%%%%%%%%%%%%%%%%%%%%%%%%%%%%%%%%%%%%%%%%%%%%
%%%%%%%%%%% BEGIN PROBLEM P2 %%%%%%%%%%%%%%%%%%%%%%%%
%%%%%%%%%%%%%%%%%%%%%%%%%%%%%%%%%%%%%%%%%%%%%%%%%%%%%

\item[\textbf{P2.}]
Let $X$ be a continuous random variable with the following probability density function.

\[p(x)=\begin{cases}cx^2\qquad -1\le x\le 2\\ 0 \qquad \text{otherwise}\end{cases}\]

\begin{enumerate}
    \item What is the support of $X$?

    \vfill

    \item Find the value of $c$ that makes $p(x)$ a valid probability density function.

    \vfill
    
    \item Find the formula of the cumulative distribution function $C(x)$ of $X$ on the support of $X$.

    \vfill

    \item Find $P(0<X<2)$.

    \vfill
\end{enumerate}

%%%%%%%%%%%%%%%%%%%%%%%%%%%%%%%%%%%%%%%%%%%%%%%%%%%%%
%%%%%%%%%%% END PROBLEM P2 %%%%%%%%%%%%%%%%%%%%%%%%%%
%%%%%%%%%%%%%%%%%%%%%%%%%%%%%%%%%%%%%%%%%%%%%%%%%%%%%
\end{enumerate}

\end{document}